\todo{Lecture 12}

\chapter{Produits scalaires et hermitiens}

\begin{definition}[Produit scalaire]
Soit $V$ un espace vectoriel sur $\mathbf{R}$ muni d'une forme bilineaire symetrique $\langle  \rangle$. $\langle  \rangle$ est definie positive si pour tout $v\in V\setminus \{ 0 \}$ on a que $\langle v,v \rangle>0$. Dans ce cas on appelle $\langle  \rangle$ un produit scalaire.
\end{definition}
\begin{exercise}
Soit $V$ un espace vectoriel sur $\mathbf{R}$ et $\langle  \rangle$ une forme bilineaire symetrique. Supposons que $u,v\in V$ et $\langle u,u \rangle>0$, $\langle v,v \rangle<0$ et $\{ u,v \}$ est libre. Trouver $x \in V\setminus \{ 0 \}$ tel que $\langle x,x \rangle=0$. 
\end{exercise}
\begin{definition}[Norme]
Soit $v\in V$. On definit la norme de $v$ comme 
$$
\lVert v \rVert =\sqrt{ \langle v,v \rangle  }
$$
\end{definition}
\begin{definition}
On dit que $v\in V$ est unitaire si $\lVert v \rVert=1$.
\end{definition}
\begin{proposition}
Soit $v\in V$ et $\alpha \in \mathbf{R}$. Alors $\lVert \alpha v \rVert=|\alpha|\lVert v \rVert$.
\end{proposition}
\begin{proof}
On a 
\begin{align*}
\lVert \alpha v \rVert^{2} & =\langle \alpha v,\alpha v \rangle \\
 & =\alpha^{2}\langle v,v \rangle \\
\end{align*}
\end{proof}
\begin{example}
Soit $V=\mathbf{R}^{n}$, $x,y\in \mathbf{R}^{n}$. $\langle x,y \rangle=\sum_{i=1}^{n}x_{i}y_{i}$ est le produit scalaire standard. Si $\dim(V)=n$ notre theorie (algorithme) nous qu'il existe une base $B=\{ b_{1},\dots,b_{n} \}$ tel que $A_{B}^{\langle  \rangle}=I_{n}$. Donc $\langle u,v \rangle=[u]_{B}^{\mathsf{T}}[v]_{B}$.
\end{example}

\begin{theorem}[Pythagore]
Soit $u,v\in V$ perpendiculaires, alors $\lVert u+v \rVert^{2}=\lVert u \rVert^{2}+\lVert v \rVert^{2}$.
\end{theorem}
\begin{proof}
\begin{align*}
\lVert u+ v\rVert ^{2} & =\langle u+v,u+v \rangle   \\
 & =\langle u,u\rangle +2\langle u,v \rangle +\langle v,v \rangle  \\
  & =\lVert u \rVert ^{2}+\lVert v \rVert ^{2}  \\
\end{align*}
\end{proof}
\begin{theorem}[Cauchy-Schwarz]
Pour tous $u,v\in V$ :
$$
	|\langle u,v \rangle |\le \lVert  u \rVert\lVert v \rVert  
$$
\end{theorem}
\begin{proof}
Si $u=0$ alors $\checkmark$. Sinon on pose $\alpha=\langle v,u \rangle/\langle u,u \rangle$ donc 
\begin{align*}
\lVert v \rVert ^{2} & =\lVert v-\alpha u  \rVert ^{2}+\alpha^{2}\lVert u \rVert ^{2} \\
 & \ge \alpha^{2}\lVert u \rVert ^{2}= \frac{\langle v,u  \rangle^{2} }{\langle u,u \rangle ^{2}} \lVert u \rVert ^{2} \\
 & =\frac{\langle v,u \rangle ^{2}}{\lVert u \rVert ^{2}}
\end{align*}
\end{proof}

\begin{theorem}
Soient $a_{1},\dots,a_{k}\in V\setminus \{ 0 \}$ deux-a-deux orthogonaux et $v\in V$. Ils existent $\beta_{1},\dots,\beta _{k}\in \mathbf{R}$ uniques, tels que 
$$
\left\langle  v-\sum_{i=1}^{k}\beta _{i}a_{i},a_{j}   \right\rangle=0,\quad \forall j\in \{ 1,\dots,k \} 
$$
\end{theorem}
\begin{proof}
\begin{align*}
\left\langle  v-\sum_{i=1}^{k}\beta _{i}a_{i},a_{j}   \right\rangle & =0 \\
 & =\langle v,a_{j} \rangle -\sum _{i=1}^{k}\beta _{i}\left\langle  a_{i} ,a_{j} \right\rangle \\
 & =\langle v,a_{j} \rangle -\beta _{j}\langle a_{j},a_{j} \rangle  \\
\Leftrightarrow \beta _{j} & = \frac{\langle v,a_{j} \rangle }{\langle a_{j},a_{j} \rangle }
\end{align*}
on appelle $\beta _{j}$ un coefficient de Fourier relativement a $a_{j}$.
\end{proof}
\begin{theorem}[Procede de Gram-Schmidt]
Soit $\{ v_{1},\dots,v_{n} \}\subset V$ un ensemble libre. Il existe un ensemble libre $\{ u_{1},\dots,u_{n} \}$ orthogonal tel que pour tout $i\in \{ 1,..,n \}$ on a que $\mathrm{span}(u_{1},\dots,u_{i})=\mathrm{span}(v_{1},\dots,v_{i})$. 
\end{theorem}
\begin{proof}
On procede par recurrence. Pour le cas $i=1$ on a $u_{1}=v_{1}$.

Pour le cas $i>1$, on suppose le theoreme vrai jusqu'a $i-1$ donc $\mathrm{span}(u_{1},\dots,u_{i-1})=\mathrm{span}(v_{1},\dots,v_{i-1})$. Alors
$$
u_{i}=v_{i}-\sum_{j=1}^{i-1} \frac{\langle v_{i},u_{j} \rangle }{\langle u_{j},u_{j} \rangle }u_{j} 
$$
donc par recurrence
$$
\mathrm{span}(v_{1},\dots,v_{i})=\mathrm{span}(u_{1},\dots,u_{i-1},v_{i})=\mathrm{span}(u_{1},\dots,u_{i-1},u_{i})
$$
\end{proof}


